\documentclass[twocolumn]{aastex631}
\begin{document}
\title{A residual reionization ionizing-photon-budget shortfall for star-forming galaxies at z\~{}7 (crisis systematic corner)}
\author{NebulaMind Lab (autonomous pipeline)}
\affiliation{NebulaMind Lab, \url{https://lab.nebulamind.net}}
\begin{abstract}
Reionization ionizing-photon-budget reconciliation at z\~{}7: star-forming galaxies require f\_esc=0.394 (+0.314/-0.174) to close the budget (Madau-Dickinson SFRD, log xi\_ion=25.5+/-0.15, clumping C=2-5, JWST-SFRD tail), versus indirect-proxy-inferred f\_esc=0.050 (+0.075/-0.030) from LzLCS O32/beta calibrations. Median delta(required-inferred)=+0.326 dex-frac (16-84\%: +0.142 to +0.643); 97\% of the systematic MC shows a shortfall. Conclusion: a genuine SHORTFALL remains, and the sign holds under both O32 and beta calibrations. The result is bounded by the xi\_ion x clumping x proxy-calibration systematic, not by statistics. Generated autonomously from public data (jwst) via the NebulaMind Lab runner: a bounded, reproducible, descriptive study.
\end{abstract}
\section{Introduction}
The ionizing photon budget during reionization has been a topic of significant interest in recent years, with studies suggesting that there may be a shortfall in the number of photons required to sustain reionization [Muñoz2024]. This potential crisis has led researchers to question whether our current understanding of star-forming galaxies and their role in reionization is sufficient [Davies2021]. In order to address this issue, we have conducted an analysis of the ionizing photon budget at z\~{}7 using a literature-anchored approach.
\section{Data and method}
Our method involves calculating the cosmic star formation rate density (SFRD) using the Madau \& Dickinson (2014) analytic fitting function. We adopt published values for xi\_ion and the O32/beta f\_esc proxy calibrations from LzLCS [Chisholm+22, Flury+22; Simmonds+24]. By reconciling these values with the ionizing photon budget, we aim to determine whether star-forming galaxies can account for the required number of photons.
\section{Result}
Our analysis reveals that in order to close the reionization ionizing-photon-budget at z\~{}7, star-forming galaxies would require an escape fraction (f\_esc) of 0.394 (+0.314/-0.174). However, indirect-proxy-inferred f\_esc values from LzLCS O32/beta calibrations suggest a significantly lower value of 0.050 (+0.075/-0.030). This discrepancy indicates a median delta of +0.326 dex-frac (16-84\%: +0.142 to +0.643), with 97\% of our systematic Monte Carlo simulations showing a shortfall in the photon budget.
\begin{figure}\centering\includegraphics[width=\columnwidth]{result.png}\caption{A residual reionization ionizing-photon-budget shortfall for star-forming galaxies at z\~{}7 (crisis systematic corner)}\end{figure}
\section{Caveats}
It is important to note that our analysis relies on automated, single-selection, and uncalibrated measurements, which may introduce limitations and uncertainties. The accuracy of our results depends heavily on the assumptions made regarding xi\_ion, clumping factor C, and proxy calibrations. Additionally, our study does not account for potential systematic errors in the literature values used or the impact of other ionizing sources beyond star-forming galaxies. These factors highlight the need for further research and more robust measurements to fully understand the reionization photon budget.
\section*{References}
\footnotesize
[Muoz2024] Muñoz et al. (2024). Reionization after JWST: a photon budget crisis?. bibcode 2024MNRAS.535L..37M \par [Davies2021] Davies et al. (2021). The Predicament of Absorption-dominated Reionization: Increased Demands on Ionizing Sources. bibcode 2021ApJ...918L..35D \par [Park2022] Park et al. (2022). Calibrating excursion set reionization models to approximately conserve ionizing photons. bibcode 2022MNRAS.517..192P \par [Duncan2015] Duncan et al. (2015). Powering reionization: assessing the galaxy ionizing photon budget at z < 10. bibcode 2015MNRAS.451.2030D \par [Madau2017] Madau (2017). Cosmic Reionization after Planck and before JWST: An Analytic Approach. bibcode 2017ApJ...851...50M

\end{document}
